\documentclass[legal]{polinetwork}

\pnTitle{Privacy Policy}
\pnSubtitle{Informativa ai sensi dell'art.~13 del Regolamento (UE) 2016/679 («GDPR»)}

% Evita righe isolate in cima o in fondo alla pagina
\widowpenalty=10000
\clubpenalty=10000

\begin{document}

\pnMakeHeader
\vspace{-16pt}
\noindent{\small Ultimo aggiornamento: 14 ottobre 2026}\par
\vspace{10pt}

Questa informativa spiega come PoliNetwork APS (di seguito «Associazione» o «Titolare») tratta i dati personali di chi usa i suoi servizi online (di seguito «Interessato»):
\begin{itemize}
    \item il sito \href{https://polinetwork.org}{polinetwork.org};
    \item l'account PoliNetwork (\href{https://auth.polinetwork.org}{auth.polinetwork.org});
    \item la piattaforma Admin Dashboard (\href{https://admin.polinetwork.org}{admin.polinetwork.org});
    \item i gruppi e i bot Telegram gestiti dall'Associazione;
    \item gli altri servizi collegati, come \href{https://aule.polinetwork.org}{aule.polinetwork.org} e \href{https://polinet.cc}{polinet.cc}.
\end{itemize}

\section{Titolare del trattamento e contatti}
\label{sec:titolare}

Il Titolare del trattamento è PoliNetwork APS, con sede in Via Aicardo 4, 20141 Milano (MI), codice fiscale 97927490157.

Per qualsiasi domanda sui propri dati personali o per esercitare i propri diritti (Sezione~\ref{sec:diritti}), l'Interessato può scrivere a \href{mailto:privacy@polinetwork.org}{privacy@polinetwork.org} oppure alla PEC \href{mailto:polinetwork@pec.it}{polinetwork@pec.it}.

\section{Dati trattati}
\label{sec:dati}

\subsection{Dati dell'account e contenuti}
\label{sec:dati-forniti}

Chi crea e usa un account PoliNetwork fornisce:
\begin{itemize}
    \item nome, indirizzo email e immagine del profilo, ricevuti da Google o dall'account Microsoft rilasciato da PoliNetwork al momento dell'accesso;
    \item se lo desidera: l'account Telegram collegato (identificativo e username), le passkey registrate e l'indirizzo email istituzionale del Politecnico di Milano, usato per verificare lo status di studente. Delle passkey viene salvata solo la chiave pubblica: impronta digitale, volto e PIN restano sul dispositivo;
    \item i contenuti che pubblica (ad es.~appunti, guide, annunci e recensioni) e i messaggi o le segnalazioni che invia al Titolare.
\end{itemize}

All'account sono inoltre associati l'eventuale appartenenza all'Associazione (socio o membro del Consiglio Direttivo) e i ruoli e permessi assegnati.

I contenuti pubblicati sono visibili agli altri utenti o a chiunque, a seconda del servizio, eventualmente insieme al nome dell'autore. Quando il servizio lo prevede, è l'Interessato a scegliere se mostrare il proprio nome.

\subsection{Dati dei gruppi Telegram}
\label{sec:dati-telegram}

Per moderare i gruppi Telegram dell'Associazione, i bot trattano i dati del profilo Telegram (identificativo, nome, username e lingua impostata), i messaggi inviati nei gruppi, le segnalazioni e i provvedimenti di moderazione.

\subsection{Dati raccolti automaticamente}
\label{sec:dati-automatici}

Durante l'uso dei servizi, i sistemi del Titolare registrano automaticamente:
\begin{itemize}
    \item indirizzo IP, tipo di browser e dispositivo (\textit{user-agent}), data e ora e pagine richieste;
    \item le sessioni di accesso, con il relativo indirizzo IP e \textit{user-agent};
    \item le operazioni svolte sulle piattaforme, in particolare quelle amministrative (chi le ha svolte e quando);
    \item l'uso dei comandi dei bot Telegram (quale comando, chi lo ha usato e in quale chat);
    \item log tecnici e di sicurezza, come errori e tentativi di accesso non autorizzati.
\end{itemize}

Il Titolare non chiede password, dati di pagamento né categorie particolari di dati personali (art.~9 GDPR, ad es.~dati sulla salute o sulle opinioni politiche).

\section{Finalità e base giuridica}
\label{sec:finalita}

Ogni trattamento si basa su una delle «basi giuridiche» previste dall'art.~6, par.~1, GDPR:
\begin{itemize}
    \item \textbf{contratto (lett.~b):} il trattamento serve a fornire il servizio richiesto dall'Interessato;
    \item \textbf{obbligo legale (lett.~c):} il trattamento è richiesto dalla legge;
    \item \textbf{legittimo interesse (lett.~f):} il trattamento tutela un interesse concreto del Titolare, come la sicurezza dei servizi, senza che prevalgano i diritti e le libertà dell'Interessato.
\end{itemize}

Le finalità e le relative basi giuridiche sono:
\begin{itemize}
    \item creare e gestire l'account, fornire i servizi richiesti e verificare lo status di studente: \textit{contratto (b)};
    \item gestire l'appartenenza all'Associazione e assegnare ruoli e permessi: \textit{contratto (b)} e \textit{legittimo interesse (f)} a impedire accessi non autorizzati;
    \item proteggere i sistemi e prevenire gli abusi: \textit{legittimo interesse (f)};
    \item moderare i contenuti e i gruppi Telegram, anche contro lo spam (Sezione~\ref{sec:decisioni-automatizzate}): \textit{legittimo interesse (f)};
    \item rispettare gli obblighi di legge e le richieste delle autorità: \textit{obbligo legale (c)};
    \item difendere i diritti dell'Associazione, anche in tribunale: \textit{legittimo interesse (f)}.
\end{itemize}

I dati non vengono usati per marketing o profilazione e non vengono venduti.

Nessuna legge obbliga l'Interessato a fornire i propri dati, ma senza i dati dell'account non è possibile usare i servizi che lo richiedono, e senza i dati raccolti automaticamente i servizi non possono funzionare. Collegare l'account Telegram, verificare lo status di studente e registrare passkey è facoltativo: chi non lo fa non potrà usare solo le funzionalità che lo richiedono.

\section{Decisioni automatizzate}
\label{sec:decisioni-automatizzate}

Per contrastare lo spam nei gruppi Telegram, i bot applicano automaticamente regole fisse, senza intelligenza artificiale né profilazione. Ad esempio, eliminano messaggi con link non consentiti o inviati in massa, silenziano gli utenti senza username Telegram, silenziano o escludono da tutti i gruppi gli account che inviano spam e rifiutano richieste di ingresso sospette.

Chi ritiene di essere stato colpito per errore può chiedere che il provvedimento venga riesaminato da un membro dello staff (Sezione~\ref{sec:titolare}).

\section{Conservazione e sicurezza}
\label{sec:conservazione}

I dati vengono conservati solo per il tempo necessario e poi cancellati o resi anonimi:
\begin{itemize}
    \item \textbf{account:} finché l'account esiste;
    \item \textbf{account collegati e status di studente:} finché restano collegati; la verifica dello status di studente vale 365 giorni;
    \item \textbf{sessioni di accesso:} scadono dopo 7 giorni senza utilizzo e vengono cancellate entro 30 giorni dalla scadenza;
    \item \textbf{messaggi nei gruppi Telegram:} 7 giorni;
    \item \textbf{dati per riconoscere lo spam su Telegram:} 30 giorni;
    \item \textbf{provvedimenti di moderazione:} 24 mesi dopo la scadenza se temporanei, fino alla revoca se permanenti;
    \item \textbf{registro delle operazioni amministrative:} 5 anni;
    \item \textbf{uso dei comandi dei bot Telegram:} 90 giorni;
    \item \textbf{log tecnici:} 30 giorni;
    \item \textbf{conteggio delle richieste per indirizzo IP:} 1 ora;
    \item \textbf{backup:} 90 giorni; i dati cancellati restano nei backup fino a questa scadenza.
\end{itemize}

I dati possono essere conservati più a lungo solo se serve per rispettare un obbligo di legge o per difendere i diritti dell'Associazione.

Il Titolare protegge i dati con misure tecniche e organizzative adeguate (art.~32 GDPR), come connessioni cifrate, messaggi e backup conservati in forma cifrata e accessi limitati in base al ruolo.

\section{Destinatari dei dati}
\label{sec:destinatari}

I dati possono essere comunicati solo a:
\begin{itemize}
    \item \textbf{volontari dell'Associazione} autorizzati (ad es.~amministratori e moderatori), per quanto serve al loro ruolo;
    \item \textbf{fornitori tecnici}, che agiscono come responsabili del trattamento (art.~28 GDPR): Microsoft (server e backup nell'Unione europea, invio delle email, account dell'Associazione), Cloudflare (rete e protezione dei siti) e GitHub (pubblicazione di \href{https://aule.polinetwork.org}{aule.polinetwork.org});
    \item \textbf{titolari autonomi}, secondo le proprie informative privacy: Google (accesso con account Google), Telegram (account collegato e gruppi) e OpenFreeMap (mappe di \href{https://aule.polinetwork.org}{aule.polinetwork.org});
    \item \textbf{autorità pubbliche}, quando previsto dalla legge.
\end{itemize}

I dati non vengono diffusi, tranne i contenuti che l'Interessato sceglie di pubblicare.

I server dell'Associazione si trovano nell'Unione europea. Alcuni fornitori (Microsoft, GitHub, Cloudflare e Google) hanno sede negli Stati Uniti: in questi casi il trasferimento si basa sul Data Privacy Framework UE-USA (decisione di adeguatezza del 10 luglio 2023, art.~45 GDPR; \href{https://www.dataprivacyframework.gov}{dataprivacyframework.gov}) e sulle clausole contrattuali standard approvate dalla Commissione europea (art.~46 GDPR). Una copia di queste garanzie può essere richiesta ai contatti indicati nella Sezione~\ref{sec:titolare}.

\section{Cookie}
\label{sec:cookie}

I servizi usano solo cookie tecnici, necessari per accedere all'account e mantenere attiva la sessione. Non ci sono cookie di profilazione né strumenti di analisi o di tracciamento. Per i cookie tecnici non serve il consenso (art.~122 D.Lgs.~196/2003), quindi non viene mostrato alcun banner.

Alcune preferenze, come il tema, la lingua e le aule preferite, restano salvate solo nel browser dell'Interessato. Bloccando i cookie non è possibile accedere ai servizi che richiedono un account.

\section{Diritti dell'Interessato}
\label{sec:diritti}

Scrivendo ai contatti indicati nella Sezione~\ref{sec:titolare}, l'Interessato può in qualsiasi momento:
\begin{itemize}
    \item accedere ai propri dati e riceverne una copia (art.~15 GDPR);
    \item correggerli o completarli (art.~16);
    \item chiederne la cancellazione, compresa quella dell'account (art.~17);
    \item chiedere che vengano solo conservati e non usati in altro modo, nei casi previsti dalla legge (art.~18);
    \item riceverli in un formato leggibile da un computer o farli inviare a un altro titolare (art.~20);
    \item chiedere che una persona riesamini le decisioni automatiche dei bot e contestarle (art.~22).
\end{itemize}

\textbf{Diritto di opposizione.} L'Interessato può opporsi in qualsiasi momento, per motivi legati alla sua situazione particolare, ai trattamenti basati sul legittimo interesse (art.~21 GDPR). In questo caso il Titolare smette di trattare i dati, a meno che non dimostri motivi legittimi prevalenti o la necessità di difendere i propri diritti.

Dalla pagina dell'account su \href{https://auth.polinetwork.org}{auth.polinetwork.org} è anche possibile chiudere le sessioni attive e scollegare gli account collegati.

Il Titolare risponde entro un mese, prorogabile di altri due mesi per le richieste complesse (art.~12 GDPR). L'Interessato può inoltre proporre reclamo al Garante per la protezione dei dati personali (\href{https://www.garanteprivacy.it}{garanteprivacy.it}).

\section{Minori}
\label{sec:minori}

Per creare un account PoliNetwork bisogna avere almeno 14 anni. Se il Titolare scopre di aver raccolto dati di una persona più giovane, li cancella.

\section{Modifiche all'informativa}
\label{sec:modifiche}

L'informativa può essere aggiornata, ad esempio quando cambiano i servizi. La versione in vigore, con la data dell'ultimo aggiornamento, è sempre disponibile su \href{https://polinetwork.org/privacy}{polinetwork.org/privacy}. In caso di modifiche importanti, gli utenti vengono avvisati tramite i servizi dell'Associazione.

\end{document}
